\documentclass[11pt,a4paper]{article}

\newcommand{\rapporttitre}{Industrialisation}
\newcommand{\rapportversion}{v0.5}
% Mise en page commune à tous les rapports du projet.
% Chaque rapport définit \rapporttitre et \rapportversion avant de charger ce fichier.

\usepackage{fontspec}
\usepackage[french]{babel}
\usepackage[a4paper,margin=2.2cm,bottom=2.6cm]{geometry}
\usepackage{graphicx}
\usepackage{float}
\usepackage{booktabs}
\usepackage{tabularx}
\usepackage{array}
\usepackage[table]{xcolor}
\usepackage{caption}
\usepackage{enumitem}
\usepackage{titlesec}
\usepackage{fancyhdr}
\usepackage{hyperref}
\usepackage{xurl}

\setmainfont{texgyreheros}[
  Extension=.otf,
  UprightFont=*-regular,
  BoldFont=*-bold,
  ItalicFont=*-italic,
  BoldItalicFont=*-bolditalic]
\definecolor{accent}{HTML}{2A78D6}
\definecolor{ink2}{HTML}{52514E}
\definecolor{fond}{HTML}{EEF4FC}
\hypersetup{colorlinks=true, linkcolor=accent, urlcolor=accent}
\graphicspath{{figures/}}

\titleformat{\section}{\Large\bfseries\color{accent}}{\thesection.}{0.6em}{}[\vspace{-0.4em}\color{accent}\rule{\linewidth}{0.8pt}]
\titleformat{\subsection}{\large\bfseries}{\thesubsection}{0.6em}{}
\titlespacing*{\section}{0pt}{1.6em}{0.9em}

\captionsetup{font={small,color=ink2}, labelfont=bf, justification=raggedright, singlelinecheck=false}
\setlist{itemsep=0.25em, topsep=0.4em}
\setlength{\parskip}{0.45em}
\renewcommand{\arraystretch}{1.25}
\newcolumntype{Y}{>{\raggedright\arraybackslash}X}

\pagestyle{fancy}
\fancyhf{}
\fancyhead[L]{\small\color{ink2} complaint-triage -- \rapporttitre}
\fancyhead[R]{\small\color{ink2} \rapportversion}
\fancyfoot[R]{\small\color{ink2}\thepage}
\renewcommand{\headrulewidth}{0.3pt}

% \pagedegarde{sous-titre}{date}{ligne de version}
\newcommand{\pagedegarde}[3]{%
  \hypersetup{pageanchor=false}
  \begin{titlepage}
    \setlength{\parindent}{0pt}
    \vspace*{5cm}
    {\color{accent}\bfseries\large COMPLAINT-TRIAGE}\par\vspace{0.6cm}
    {\fontsize{30}{36}\selectfont\bfseries \rapporttitre}\par\vspace{0.5cm}
    {\Large\color{ink2} #1}\par
    \vfill
    {\color{accent}\rule{\linewidth}{1.2pt}}\par\vspace{0.3cm}
    {\color{ink2} Nathan Mang\\ #2\\ #3}
    \vspace{2cm}
  \end{titlepage}
  \hypersetup{pageanchor=true}}

% Encadré pour les points à retenir.
\newcommand{\aretenir}[1]{%
  \par\noindent\colorbox{fond}{\parbox{\dimexpr\linewidth-2\fboxsep}{\small #1}}\par}

\usepackage{microtype}

\begin{document}

\pagedegarde{La base et l'image de l'application testées à chaque modification}{Octobre 2026}{Rapport de fin de version}

\section*{Synthèse}
\addcontentsline{toc}{section}{Synthèse}

Jusqu'à la v0.4, les tests d'intégration ne tournaient que sur le poste de développement, où PostgreSQL est installé, et l'application n'existait pas sous forme d'image. Cette version, volontairement légère, fait tourner les deux dans la CI de GitHub à chaque PR, sans rien installer sur le poste et sans aucun coût. La CI compte désormais trois tâches en parallèle, toutes vertes dès le premier lancement : les contrôles et les tests unitaires avec les scénarios de sûreté (66~secondes), une base PostgreSQL jetable avec les tests d'intégration (37~secondes) et l'image Docker de l'application (63~secondes).

\setcounter{tocdepth}{1}
\tableofcontents
\newpage

\section{Objectif et choix}

\subsection{Une version légère, par choix}

Le cadrage prévoyait pour la v0.5 Docker et une chaîne de déploiement complète. Deux contraintes ont conduit à une version plus légère. Le poste de développement n'a que 8~Go de mémoire, et Docker Desktop en occupe à lui seul plusieurs. Et le projet reste sur une règle de zéro dépense : un déploiement sur Google Cloud exige d'y lier un moyen de paiement. Le but de la version est donc ciblé : que la base et l'image de l'application soient testées automatiquement, dans la CI de GitHub, gratuite pour un dépôt public.

\subsection{Ce qui est fait, ce qui est reporté}

\begin{table}[H]
\centering\small
\begin{tabularx}{\linewidth}{@{}Y Y@{}}
\toprule
Fait & Reporté, ou écarté \\
\midrule
PostgreSQL jetable dans la CI, mis en place par le script du projet ; tests d'intégration à chaque PR & Docker sur le poste : inutile tant que la CI construit et teste l'image \\
Image Docker de l'application, construite et vérifiée à chaque PR & Publication de l'image dans un registre : utile seulement avec un hébergement \\
Versions exactes de Python, uv et PostgreSQL & Déploiement sur Google Cloud : remplacé en v1.0 par une démo publique gratuite, sans moyen de paiement \\
\bottomrule
\end{tabularx}
\caption{Périmètre de la v0.5}
\end{table}

\section{La base dans la CI}

À chaque PR, GitHub démarre un conteneur PostgreSQL 18.6, la même version que sur le poste, et attend qu'il réponde. Le script de mise en place du projet crée ensuite les comptes, le schéma et les droits, exactement comme sur le poste. C'est aussi un test du script lui-même, sur une machine vierge. Sur le poste, le mot de passe du super-utilisateur reste demandé en saisie masquée ; en CI, il vient d'une variable d'environnement. La base est détruite à la fin de la tâche, et son mot de passe n'est donc pas un secret.

Les dix-huit tests d'intégration tournent alors : les droits de chaque compte, vérifiés en tentant les opérations interdites, et les outils de l'agent sur une banque fictive. Un piège demandait une précaution. Sans base, ces tests sont sautés, et un test sauté ne fait pas échouer la CI : une configuration défaillante passerait donc au vert sans rien avoir testé. Dans cette tâche, une variable \texttt{REQUIRE\_DATABASE} empêche ce saut ; un test privé de base échoue.

\section{L'image Docker}

L'image contient le paquet, les scripts et les jeux d'évaluation. Ses choix suivent quatre principes :
\begin{itemize}
  \item \textbf{des versions exactes} : Python 3.12.15, la plus récente au moment de la version, et uv 0.13.1, comme sur le poste et dans la CI ; toutes deux ont été vérifiées sur les registres avant d'être figées ;
  \item \textbf{les dépendances du fichier \texttt{uv.lock}}, installées dans une couche à part, gardée en cache tant que ce fichier ne change pas ;
  \item \textbf{une exécution sans root}, par un utilisateur dédié ;
  \item \textbf{rien de privé} : un fichier \texttt{.dockerignore} écarte le fichier \texttt{.env}, les données, les environnements, les documents et les tests. Les données se montent au lancement, les secrets passent par les variables d'environnement.
\end{itemize}

La CI construit l'image, puis vérifie qu'elle démarre, qu'elle importe l'agent et le RAG, qu'elle ne tourne pas en root, et qu'elle ne contient ni fichier \texttt{.env}, ni données, ni tests.

\section{Bilan}

\begin{table}[H]
\centering\small
\begin{tabularx}{\linewidth}{@{}l Y r@{}}
\toprule
Tâche & Contenu & Durée \\
\midrule
Lint, types et tests & ruff, mypy, 205 tests unitaires, 17 scénarios de sûreté & 66~s \\
Base et tests d'intégration & PostgreSQL 18.6, script de mise en place, 18 tests & 37~s \\
Image Docker & Construction et quatre vérifications & 63~s \\
\bottomrule
\end{tabularx}
\caption{La CI à la fin de la v0.5, premier lancement}
\end{table}

Chaque modification du code passe désormais par les mêmes contrôles qu'un projet d'équipe : qualité, sûreté de l'agent, droits de la base, et image de l'application. La v1.0 utilisera cette image pour une démo publique gratuite.

\end{document}
